\subsection{The two common scales}

\begin{lemma}[Material utility and a common affine scale]\label{lf:lem:material-scale}
Under Axioms \textup{(A1)--(A3)} and
\textup{(A5)--(A8)}, there are outcomes $o^+,o^-$ and a
nonconstant function $u:O\to\R$ such that
\[
u(o^+)=1,
\qquad
u(o^-)=0.
\]
For every payoff lottery $\ell\in\D(O)$, the payoff-lottery order is represented
by $\sum_o\ell(o)u(o)$.  For every full-support prior $q$, the representatives
in Lemma~\ref{lf:lem:terminal-affine} may be normalised so that
\begin{equation}
W_q(\ell)=\sum_o\ell(o)u(o).
\label{lf:eq:payoff-normalisation}
\end{equation}
With this normalisation, for every full-support $q\in\D(A)$, every nonempty
finite $B$, every full-support $s\in\D(B)$, and every finite experiment $E$ at
$q$,
\[
W_{q\otimes s}(E^\uparrow)=W_q(E),
\]
where $E^\uparrow$ ignores the dummy coordinate.  Moreover, for every pair of nonempty
finite action sets, every two full-support priors $q,q'$ on those sets, and
every two finite experiments $E,G$ on the respective sets,
\begin{equation}
(q,E)\succeq(q',G)
\quad\Longleftrightarrow\quad
W_q(E)\ge W_{q'}(G).
\label{lf:eq:common-scale}
\end{equation}
\end{lemma}

\begin{proof}
On a singleton action set, terminal-law sufficiency identifies an experiment
with its payoff lottery.  Public-mixture independence and continuity, proved as
in Lemma~\ref{lf:lem:terminal-affine}, give an affine utility on $\D(O)$.
By Lemma~\ref{lem:relevance-bridge}, Axiom~\textup{(A3)} supplies
$o^+\succ o^-$.  Normalise their values to one
and zero and let $u(o)$ be the value of the degenerate lottery at $o$.
Affinity on the finite simplex gives expected utility.

An action processor may ignore its input and draw an arbitrary target prior.
Axiom~\textup{(A7)}, applied in both directions, therefore makes this payoff
lottery order independent of the prior and action alphabet.  Normalise each
$W_q$ at the same two deterministic payoffs.  The payoff-lottery embedding is
affine and order-reflecting, so positive-affine uniqueness gives
\eqref{lf:eq:payoff-normalisation}.

Dummy lifting between full-support fibres is likewise an affine order
embedding.  Positive-affine uniqueness initially permits
\[
W_{q\otimes s}(E^\uparrow)=cW_q(E)+d,
\qquad c>0.
\]
The two payoff anchors have values zero and one on both sides; hence $c=1$ and
$d=0$.  To compare $E$ at $q$ with $G$ at $q'$, lift each to the common product
prior $q\otimes q'$, using the other prior as the dummy.  Equation
\eqref{lf:eq:dummy-ordinal}, the replacement part of
Lemma~\ref{lf:lem:bookkeeping}, and the common product representative give
\eqref{lf:eq:common-scale}.
\end{proof}

Retain the outcome $o_\ast$ fixed in the preceding pure-trace subsection to
witness Axiom~\textup{(A4)}.  Proposition~\ref{pt:prop:main} has established
that mutual information represents every within-channel and block-supported
comparison on this constant-payoff face.
The normalization $u(o^+)=1$ and $u(o^-)=0$ fixes only the material origin
and unit.  The trace anchor $o_\ast$ need not be either material anchor, and
its intercept $u(o_\ast)$ is retained.  It remains to identify the cardinal
coefficient of mutual information relative to the material unit and to show
that this coefficient is common to all priors and finite alphabets.

\begin{lemma}[The trace scale]\label{lf:lem:trace-scale}
Under Axioms \textup{(A1)--(A8)}, there is one number $\lambda>0$, independent
of the prior and all finite alphabets, such that, for every nonempty finite
$A$ with $\lvert A\rvert\ge2$, every full-support $q\in\D(A)$, every nonempty
finite record set $R$, and every experiment
$E:A\to\D(O\times R)$ that pays $o_\ast$ surely,
\begin{equation}
W_q(E)=u(o_\ast)+\lambda I_q(E).
\label{lf:eq:sure-anchor}
\end{equation}
\end{lemma}

\begin{proof}
Fix a nontrivial full-support $q$.  On the constant-$o_\ast$ face, $W_q$ is
affine by Lemma~\ref{lf:lem:terminal-affine}.  For any experiment $E$, the
finite entropy identity
\begin{equation}
I_q(E)=\Sh(q)-\sum_{(o,\pi)\in\supp\eta_E}
\eta_E(o,\pi)\Sh(\pi).
\label{lf:eq:posterior-entropy}
\end{equation}
shows that mutual information depends only on the marked terminal law and is
affine under public mixtures.  Proposition~\ref{pt:prop:main} says that $W_q$ and
$I_q$ represent the same nonconstant order on the constant-$o_\ast$ face.
Positive-affine uniqueness therefore gives
\[
W_q(E)=\lambda_q I_q(E)+b_q,
\qquad \lambda_q>0.
\]
The uninformative experiment has mutual information zero and, by
\eqref{lf:eq:payoff-normalisation}, value $u(o_\ast)$.  Thus $b_q=u(o_\ast)$.

The coefficient does not depend on $q$.  Let $q$ and $q'$ be nontrivial
full-support priors.  Lift full revelation at $q$ to $q\otimes q'$ by adding an
independent dummy coordinate.  Normalised value and mutual information are both
unchanged.  Since the common information value is $\Sh(q)>0$, cancellation
gives $\lambda_{q\otimes q'}=\lambda_q$.  Lifting full revelation from the
other coordinate gives $\lambda_{q\otimes q'}=\lambda_{q'}$.  Hence all such
coefficients agree.  Equivalently, fix the uniform prior on a two-point dummy
alphabet and call its coefficient $\lambda$.  Axiom~\textup{(A4)} fixes the
strict orientation, so $\lambda>0$.

Finally, any experiment that pays $o_\ast$ surely is a pure record experiment
after the constant payoff coordinate is deleted; restoring that coordinate
changes neither its marked terminal law nor its mutual information.  This
proves \eqref{lf:eq:sure-anchor}.
\end{proof}

\subsection{Exact branch weights}

Let $A$ and $Y$ be arbitrary nonempty finite sets, let $q\in\D(A)$ have full
support with $\lvert A\rvert\ge2$, and let
$P:A\to\D(Y)$ be a first-stage record channel.  For every $y\in Y$, put
\[
m_y:=\sum_{a\in A}q(a)P(y\mid a).
\]
Whenever $m_y>0$, define the posterior on $A$ and its support by
\[
\widehat r_y(a):=\frac{q(a)P(y\mid a)}{m_y},
\qquad
S_y:=\supp\widehat r_y,
\]
and write $r_y\in\D(S_y)$ for the full-support restriction of
$\widehat r_y$ to $S_y$.

We need notation for changing a continuation without fixing what happens on
the other branches.  For a finite continuation experiment
$E:S_y\to\D(O\times R_E)$, choose a retraction
$\rho_y:A\to S_y$ whose restriction to $S_y$ is the identity and set
$E^\uparrow(\cdot\mid a):=E(\cdot\mid\rho_y(a))$.  Given an arbitrary
background $\mathbf H=(H_z)_{z\ne y}$ of continuation channels on $A$, let
\[
K_z^E:=
\begin{cases}
E^\uparrow,&z=y,\\
H_z,&z\ne y,
\end{cases}
\qquad
J_{y,\mathbf H}(E):=P\triangleright(K_z^E)_{z\in Y}.
\]
Here and below, fix any ordering of $Y$ when writing the compound.
Before taking the compound, all continuation record alphabets are embedded as
tagged copies in one finite disjoint union.  Payoff-preserving record
isomorphisms make the choice of tags immaterial.  Also,
$a\notin S_y$ implies $P(y\mid a)=0$, because $q$ has full support; hence the
rows of $E^\uparrow$ outside $S_y$ are never used on branch $y$.
The compound channel itself is therefore independent of the chosen
retraction.

\begin{lemma}[Exact reached-branch weight]\label{lf:lem:branch-weight}
Under Axioms \textup{(A1)--(A8)}, for every nonempty finite $A,Y$, every
full-support $q\in\D(A)$ with $\lvert A\rvert\ge2$, every
$P:A\to\D(Y)$, every reached $y\in Y$, every two finite continuation
experiments $E,G$ on $S_y$, and every finite background
$\mathbf H=(H_z)_{z\ne y}$,
\begin{equation}
W_q\bigl(J_{y,\mathbf H}(E)\bigr)
-W_q\bigl(J_{y,\mathbf H}(G)\bigr)
=m_y\bigl[W_{r_y}(E)-W_{r_y}(G)\bigr].
\label{lf:eq:branch-weight}
\end{equation}
\end{lemma}

\begin{proof}
At the branch posterior $\widehat r_y$, support restriction identifies
$(\widehat r_y,E^\uparrow)$ with $(r_y,E)$, and likewise for $G$.
Corollary~\ref{pt:cor:one-branch-a8} therefore gives, for every $E$ and $G$,
\[
E\succeq_{r_y}G
\quad\Longleftrightarrow\quad
J_{y,\mathbf H}(E)\succeq_qJ_{y,\mathbf H}(G),
\]
with the same equivalence for strict preference.  This statement includes
continuations with different record alphabets: the tagged disjoint union puts
them in the common comparison environment required by the corollary.

To see the affine structure explicitly, let
$\iota_y:\D(S_y)\to\D(A)$ pad a distribution by zeros outside $S_y$.
There is a fixed subprobability measure $\beta_{y,\mathbf H}$, of mass
$1-m_y$, containing the terminal atoms from all branches other than $y$, such
that
\[
\eta_{J_{y,\mathbf H}(E)}
=\beta_{y,\mathbf H}
 +m_y(\mathrm{id}_O,\iota_y)_\#\eta_E.
\]
Thus insertion is affine on marked terminal laws.  If $S_y$ is nontrivial,
the local order is nonconstant: at payoff $o_\ast$, full revelation and
silence have, by Lemma~\ref{lf:lem:trace-scale}, value difference
$\lambda\Sh(r_y)>0$.  Affine
uniqueness consequently gives, for every fixed background $\mathbf H$,
constants $c_{y,\mathbf H}>0$ and $d_{y,\mathbf H}\in\R$ such that
\begin{equation}
W_q\bigl(J_{y,\mathbf H}(E)\bigr)
=c_{y,\mathbf H}W_{r_y}(E)+d_{y,\mathbf H}.
\label{lf:eq:insertion-affine}
\end{equation}

It remains to identify the coefficient and to show that the background cannot
alter it.  Let $\mathbf H^\ast$ pay $o_\ast$ for sure and reveal only a
singleton record after every $z\ne y$.  For every continuation $E$, the
mutual-information chain rule gives
\begin{equation}
I_q\bigl(J_{y,\mathbf H^\ast}(E)\bigr)
=I_q(A;Y)+m_y I_{r_y}(E).
\label{lf:eq:insertion-chain}
\end{equation}
Take for $E$ full revelation, $E^{\mathrm{id}}$, and for $G$ silence,
$E^0$, both at the constant payoff $o_\ast$.  By
Lemma~\ref{lf:lem:trace-scale}, their local value difference is
$\lambda\Sh(r_y)$, while \eqref{lf:eq:insertion-chain} makes their aggregate
value difference $\lambda m_y\Sh(r_y)$.  Comparing these differences in
\eqref{lf:eq:insertion-affine}, and using
$\lambda\Sh(r_y)>0$, yields
$c_{y,\mathbf H^\ast}=m_y$.

Now fix arbitrary $E,G$, and an arbitrary background $\mathbf H$.  The crossed
public half-mixtures
\[
\tfrac12J_{y,\mathbf H}(E)
\oplus\tfrac12J_{y,\mathbf H^\ast}(G)
\quad\text{and}\quad
\tfrac12J_{y,\mathbf H}(G)
\oplus\tfrac12J_{y,\mathbf H^\ast}(E)
\]
have the same marked terminal law.  Indeed, each contains one half of each
background contribution and, on branch $y$, one half of each of
$m_y\eta_E$ and $m_y\eta_G$ (with posterior supports padded by $\iota_y$).
Lemma~\ref{lf:lem:terminal-affine}\textup{(a)} makes the mixtures indifferent.
Affinity of $W_q$ and cancellation therefore give
\[
W_q(J_{y,\mathbf H}(E))-W_q(J_{y,\mathbf H}(G))
=W_q(J_{y,\mathbf H^\ast}(E))-W_q(J_{y,\mathbf H^\ast}(G)).
\]
Applying \eqref{lf:eq:insertion-affine} to the right-hand side with
$c_{y,\mathbf H^\ast}=m_y$ proves \eqref{lf:eq:branch-weight} whenever
$\lvert S_y\rvert\ge2$.

If $S_y$ is a singleton, let $s$ be a full-support prior on a two-point set
$B$.  Adjoin $B$ independently to the global prior, lift the first-stage
channel by $P^\uparrow(y\mid a,b)=P(y\mid a)$, and lift every local and
background continuation so that it ignores $b$.  The reached posterior is
$r_y\otimes s$, its support is nontrivial, and its branch mass remains $m_y$.
The preceding argument applies to this common dummy extension.
Lemma~\ref{lf:lem:material-scale} says that adjoining or deleting the
independent dummy preserves both aggregate values and both local values.
Removing it gives \eqref{lf:eq:branch-weight} for the singleton case.
\end{proof}

For a deterministic payoff profile $g:Y\to O$, let $C(P,g)$ denote the
compound that first runs $P$ and then pays $g(y)$ with no further record.
Write $g_y^o$ for the profile that pays $o$ at $y$ and $o_\ast$ elsewhere.

\begin{lemma}[Deterministic branch assembly]
\label{lf:lem:deterministic-assembly}\label{lf:lem:branch-increment}
Under Axioms \textup{(A1)--(A8)}, for every nonempty finite $A,Y$, every
full-support $q\in\D(A)$ with $\lvert A\rvert\ge2$, every
$P:A\to\D(Y)$, and every deterministic payoff profile $g:Y\to O$,
\begin{equation}
W_q(C(P,g))
=\lambda I_q(A;Y)+\sum_{y\in Y}m_yu(g(y)).
\label{lf:eq:deterministic-assembly}
\end{equation}
\end{lemma}

\begin{proof}
Fix $y\in Y$ and $o\in O$.
If $m_y=0$, changing the payoff after $y$ leaves the marked terminal law
unchanged.  If $m_y>0$, apply Lemma~\ref{lf:lem:branch-weight} with
deterministic local payoffs $o$ and $o_\ast$ and with every other branch paying
$o_\ast$.  Equation~\eqref{lf:eq:payoff-normalisation} gives the local value
difference, and hence
\begin{equation}
W_q\bigl(C(P,g_y^o)\bigr)
-W_q\bigl(C(P,o_\ast)\bigr)
=m_y\bigl[u(o)-u(o_\ast)\bigr].
\label{lf:eq:branch-increment}
\end{equation}
The arbitrary-background clause of Lemma~\ref{lf:lem:branch-weight} gives the
same increment when the payoffs on the other branches have already been
changed.  Enumerate the reached branches and change their payoffs from
$o_\ast$ to $g(y)$ one at a time; unreached branches have no effect.  The
resulting telescope is
\begin{equation}
W_q(C(P,g))-W_q(C(P,o_\ast))
=\sum_{y\in Y}m_y\bigl[u(g(y))-u(o_\ast)\bigr].
\label{lf:eq:telescope}
\end{equation}
The baseline compound pays $o_\ast$ surely and its complete record is $Y$.
Lemma~\ref{lf:lem:trace-scale} therefore gives
$W_q(C(P,o_\ast))=u(o_\ast)+\lambda I_q(A;Y)$.
Substitution in \eqref{lf:eq:telescope}, followed by
$\sum_{y\in Y}m_y=1$, proves the formula.
\end{proof}

\subsection{Proof of the theorem}

\begin{proof}[Proof of Theorem~\ref{thm:main}]
Assume first Axioms \textup{(A1)--(A8)}.  Let $q$ have full support on a
nontrivial action set, and let $K:A\to\D(O\times R)$ be arbitrary.  Put
$Y=O\times R$ and take as first stage
\[
P_K(y\mid a):=K(y\mid a),
\qquad y=(o,r).
\]
After $y=(o,r)$, pay $o$ for sure and reveal nothing further.  Call the
resulting compound $K^{\mathrm{seq}}$.  Conditional on its final record
$((o,r),\ast)$, the posterior over actions is exactly the posterior conditional
on $(o,r)$ under $K$, and both records carry payoff $o$.  Thus $K$ and
$K^{\mathrm{seq}}$ have the same marked terminal law.  By
Lemma~\ref{lf:lem:terminal-affine}\textup{(a)},
$(q,K)\sim(q,K^{\mathrm{seq}})$.

Apply Lemma~\ref{lf:lem:deterministic-assembly} with $g(o,r)=o$.  The two
identities needed are
\[
I_q(A;Y)=I_{qK}(A;O,R)
\]
and
\[
\sum_{o,r}p_{qK}(o,r)u(o)=\E_{qK}[u(O)],
\]
respectively.  Together with sequentialisation, they give
\begin{equation}
W_q(K)=\E_{qK}[u(O)]+\lambda I_{qK}(A;O,R).
\label{lf:eq:full-support-value}
\end{equation}

If $A$ is a singleton, mutual information is zero and the marked terminal law
is that of the induced payoff lottery.  Equation
\eqref{lf:eq:payoff-normalisation} gives
\eqref{lf:eq:full-support-value} directly.  If $q$ is on the boundary, put
$S=\supp q$ and write $(q^S,K^S)$ for the restriction.  Then
\begin{equation}
(q,K)\sim(q^S,K^S),
\qquad
V(q,K)=V(q^S,K^S).
\label{lf:eq:support-value}
\end{equation}
The first identity is Lemma~\ref{lf:lem:bookkeeping}\textup{(b)}; the second
follows directly from the finite sums defining expected utility and mutual
information.  Thus no representative $W_q$ is assigned to a boundary fibre.

For arbitrary finite experiments $(q,K)$ and $(p,L)$, restrict both priors to
their supports, apply
\eqref{lf:eq:common-scale} and \eqref{lf:eq:full-support-value} there, and use
\eqref{lf:eq:support-value} to transport the comparison back.  This gives
\begin{equation}
(q,K)\succeq(p,L)
\quad\Longleftrightarrow\quad
V(q,K)\ge V(p,L),
\label{lf:eq:pair-value}
\end{equation}
with the same $u$ and $\lambda$ for every pair of finite alphabets.  The
duplication clause of Axiom~\textup{(A5)} turns this pair formula into the
within-channel representation.  Its irrelevant-block clause applies the same
formula to any two block-supported priors in any finite block environment.
This is exactly the theorem's same-scale moreover clause.  The construction
made $u$ nonconstant and $\lambda>0$, proving the forward implication.

Conversely, suppose that a nonconstant $u:O\to\R$ and $\lambda>0$ represent
every finite environment by $V$.  Completeness and transitivity of the order
on $\R$ give Axiom~\textup{(A1)}.  On every fixed finite
$(A,O,R)$, expected utility and mutual information are jointly continuous in
$(K,q)$; for mutual information, use its finite entropy formula with
$0\log0=0$.  Therefore
\[
\{(K,q,p):V(q,K)\ge V(p,K)\}
\]
is closed, which is Axiom~\textup{(A2)}.

For a duplicated channel, a prior supported on either tagged copy has the same
payoff law and information as its untagged prior.  This proves both directions
of the duplication clause in Axiom~\textup{(A5)}.  More generally, under a
prior supported on block $i$ of $\bigsqcup_{k\in I}K_k$, the block label is
constant and
\[
V\left(q_i^i,\bigsqcup_{k\in I}K_k\right)=V(q_i,K_i).
\]
The same identity for block $j$ proves both directions of the
irrelevant-block clause, for every finite $I$ and every ordered pair
$i\ne j$.  Thus Axiom~\textup{(A5)} holds.

A payoff-preserving record processor leaves the payoff marginal unchanged and,
by the Markov chain
$A-(O,R)-(O,R')$, weakly lowers mutual information.  This is
Axiom~\textup{(A6)}.  For every action processor $S:A\to\D(B)$ and every
Bayesian pushforward in \eqref{eq:action-processing}, the payoff--record
marginal is again unchanged and the undivided joint law forms the Markov chain
\[
(O,R)-A-B.
\]
Hence $I(B;O,R)\le I(A;O,R)$, proving Axiom~\textup{(A7)}.  A completion on a
processed action with $(qS)(b)=0$ has zero joint mass, so this argument covers
every completion allowed by the axiom.

For Axiom~\textup{(A8)}, fix every binary first-stage channel and continuation
triple covered by that axiom, and put
\[
\alpha:=\sum_{a\in A}q(a)P(1\mid a)>0.
\]
The expected-utility and mutual-information chain rules give the exact
difference
\begin{align*}
&V\bigl(q,P\triangleright(K,M)\bigr)
 -V\bigl(q,P\triangleright(L,M)\bigr)\\
&\qquad=\alpha\bigl[V(q_1,K)-V(q_1,L)\bigr].
\end{align*}
The first-stage information term and the common continuation after record $2$
cancel.  Since $\alpha>0$, the aggregate difference is nonnegative exactly
when the conditional difference is nonnegative, and it is positive exactly
when the conditional difference is positive.  This proves the weak
biconditional in Axiom~\textup{(A8)} and the strict biconditional.

Finally, nonconstancy of $u$ supplies distinct $o^+,o^-\in O$ with
$u(o^+)>u(o^-)$.  In the two-action channel specified by
Axiom~\textup{(A3)}, the two pure lotteries have those values and zero mutual
information, proving material relevance.  For any $o\in O$, the two lotteries
in the four-action channel specified by Axiom~\textup{(A4)} both pay $o$
surely, while their information values are $\log2$ and zero.  Because
$\lambda>0$, the required strict comparison holds (in fact at every $o$).
This verifies Axioms~\textup{(A3)} and~\textup{(A4)} and completes the
converse.  The displayed block identity also proves the theorem's moreover
clause with the same $u$ and $\lambda$.
\end{proof}
